% !TeX root = ./thuthesis.tex

% 论文基本信息配置

% 技术方案基本信息配置(简化版)

% 载入所需的宏包

% 定理类环境宏包
\usepackage{amsthm}
% 注意：不使用 amssymb，因为与 unicode-math 冲突
% 可见 unicode-math 包已经提供了所有数学符号

% 定义缺失的数学符号
% unicode-math 环境下重新定义 Box 符号
\renewcommand{\Box}{\square}

% 也可以使用 ntheorem
% \usepackage[amsmath,thmmarks,hyperref]{ntheorem}

% 注意：thuthesis.cls 已经定义了 theorem, proposition, lemma, corollary 等环境
% 所以我们不需要重新定义这些环境

% 添加 TikZ 支持
\usepackage{tikz}
\usetikzlibrary{shapes,arrows,positioning}

% 添加浮动体屏障支持
\usepackage{placeins}

% 添加代码列表支持
\usepackage{listings}
\usepackage{xcolor}

% 定义Python代码高亮颜色
% 定义现代代码高亮颜色
\definecolor{codebackground}{rgb}{0.96,0.96,0.97}
\definecolor{codecomment}{rgb}{0.5,0.5,0.5}
\definecolor{codekeyword}{rgb}{0.0,0.0,0.8}
\definecolor{codestring}{rgb}{0.0,0.5,0.0}
\definecolor{codenumber}{rgb}{0.6,0.6,0.6}
\definecolor{codeidentifier}{rgb}{0.2,0.2,0.2}

% 定义 Dart 语言 (listings 默认不支持)
\lstdefinelanguage{Dart}{
    morekeywords={abstract, as, assert, async, await, break, case, catch, class, const, continue, covariant, default, deferred, do, dynamic, else, enum, export, extends, extension, external, factory, false, final, finally, for, Function, get, hide, if, implements, import, in, interface, is, late, library, mixin, new, null, on, operator, part, rethrow, return, set, show, static, super, switch, sync, this, throw, true, try, typedef, var, void, while, with, yield},
    sensitive=true,
    morecomment=[l]{//},
    morecomment=[s]{/*}{*/},
    morestring=[b]",
    morestring=[b]'
}

% 定义通用代码样式
\lstdefinestyle{codestyle}{
    backgroundcolor=\color{codebackground},   % 背景色
    commentstyle=\color{codecomment}\itshape, % 注释样式
    keywordstyle=\color{codekeyword}\bfseries,% 关键词样式
    stringstyle=\color{codestring},           % 字符串样式
    basicstyle=\ttfamily\footnotesize\color{codeidentifier}, % 基础字体
    breakatwhitespace=false,
    breaklines=true,                 % 自动换行
    captionpos=b,                    % 标题位置
    keepspaces=true,
    numbers=left,                    % 行号位置
    numbersep=8pt,                   % 行号间距
    numberstyle=\tiny\color{codenumber}, % 行号样式
    showspaces=false,
    showstringspaces=false,
    showtabs=false,
    tabsize=2,                       % 制表符宽度
    frame=single,                    % 边框
    frameround=tttt,                 % 圆角边框
    rulecolor=\color{gray!30},       % 边框颜色
    framesep=5pt,                    % 内容与边框间距
    xleftmargin=1.5em,               % 左侧缩进
    xrightmargin=1em,                % 右侧缩进
    framexleftmargin=1em,            % 边框左侧内补
    aboveskip=1.5em,                 % 上方间距
    belowskip=1em                    % 下方间距
}

% 数学字体配置已简化，使用默认设置

% 可以使用 nomencl 生成符号和缩略语说明
% \usepackage{nomencl}
% \makenomenclature

% 表格加脚注
\usepackage{threeparttable}

% 表格中支持跨行
\usepackage{multirow}

% 固定宽度的表格。
% \usepackage{tabularx}

% 跨页表格
\usepackage{longtable}
\usepackage{tabularx}

% 算法
\usepackage{algorithm}
\usepackage{algpseudocode}

% 量和单位
% 注意: 如果使用 unicode-math,需要在 siunitx 之前加载
% \usepackage{siunitx}

% 参考文献使用 BibTeX + natbib 宏包
% 顺序编码制
\usepackage[sort]{natbib}
\bibliographystyle{thuthesis-numeric}
\usepackage{titlesec}

% 让 \paragraph 参与编号（段级别=4）
\setcounter{secnumdepth}{4}

% 定义 \paragraph 的编号样式为纯阿拉伯数字：1, 2, 3, ...
\makeatletter
\renewcommand\theparagraph{\arabic{paragraph}}
% 如果希望每个 section 内从 1 重新开始编号，保留下一行；
% 如果希望全文连续编号，请删掉下一行。
\@addtoreset{paragraph}{section}
\makeatother

% 把 \paragraph 改成块级标题（标题后换行）
\titleformat{\paragraph}[block]
{\normalfont\normalsize\bfseries}% 样式
{\theparagraph}%                  % 显示的编号（比如 1, 2, 3）
{0.8em}%                          % 编号与标题的间距
{}                                % 标题前置代码

% 调整段前段后间距（可按需修改）
\titlespacing*{\paragraph}
{0pt}{3.25ex plus 1ex minus .2ex}{1ex plus .2ex}
\usepackage{indentfirst}  % 首行缩进

% 著者-出版年制
% \usepackage{natbib}
% \bibliographystyle{thuthesis-author-year}

% 生命科学学院要求使用 Cell 参考文献格式（2023 年以前使用 author-date 格式）
% \usepackage{natbib}
% \bibliographystyle{cell}

% 本科生参考文献的著录格式
% \usepackage[sort]{natbib}
% \bibliographystyle{thuthesis-bachelor}

% 参考文献使用 BibLaTeX 宏包
% \usepackage[style=thuthesis-numeric]{biblatex}
% \usepackage[style=thuthesis-author-year]{biblatex}
% \usepackage[style=gb7714-2015]{biblatex}
% \usepackage[style=apa]{biblatex}
% \usepackage[style=mla-new]{biblatex}
% 声明 BibLaTeX 的数据库
% \addbibresource{ref/refs.bib}

% 定义所有的图片文件在 figures 子目录下
\graphicspath{{figures/}}

% 数学命令
\makeatletter
\newcommand\dif{%  % 微分符号
    \mathop{}\!%
    \ifthu@math@style@TeX
        d%
    \else
        \mathrm{d}%
    \fi
}
\makeatother

% hyperref 宏包在最后调用
\usepackage{hyperref}

% cleveref 宏包必须在 hyperref 之后加载
\usepackage[nameinlink, noabbrev]{cleveref}

% ---- 等宽字体安全回退（可选）----
% 若模板因选项/误判试图使用 Menlo 而当前系统不存在，则强制改为 Courier New
\IfFontExistsTF{Menlo}{}{\setmonofont{Courier New}[Scale=MatchLowercase]}

% =========================================
% 修复 main.tex 编译报错的补充定义
% =========================================

% 1. 定义 R 和 C 列格式 (依赖 array 宏包，已由 tabularx 加载)
\newcolumntype{R}[1]{>{\raggedleft\arraybackslash}p{#1}}
\newcolumntype{C}[1]{>{\centering\arraybackslash}p{#1}}

% 2. 定义标题字体 (使用黑体)
\newcommand{\HeadingFont}{\heiti}

% 3. 定义 1em 间距
\newcommand{\OneEm}{\hspace{1em}}

% 4. 定义带下划线的居中填空域
% 用法: \CenteredField[宽度]{内容}
\newcommand{\CenteredField}[2][10cm]{\underline{\makebox[#1][c]{#2}}}


% ---- 结束 ----